\documentclass[11pt,a4paper]{article}
\usepackage[margin=24mm]{geometry}
\usepackage{amsmath,amssymb,amsthm,booktabs,longtable,hyperref}
\setlength{\emergencystretch}{2em}
\newtheorem{theorem}{Theorem}
\newtheorem{lemma}[theorem]{Lemma}
\newcommand{\Sym}{\operatorname{Sym}}
\title{The explicit classifications at defects one to four}
\author{Human Compression project: an attributed recovery of L042}
\date{}
\begin{document}\maketitle

We recover the principal classification in Fernandes, Leemans and Mixer,
\emph{An extension of the classification of high rank regular polytopes},
arXiv:1509.01032v1, 3 September 2015. The fixed edition has 26 pages;
the subsequent journal citation is \emph{Transactions of the American
Mathematical Society} 370 (2018), 8833--8857. The explicit diagrams,
and not merely their numbers, form part of the result proved here.
The classification theorem and its original families remain attributed
to those authors.

A marking is an ordered tuple of nonidentity involutions satisfying
the string relations and the intersection property. We identify two
markings by an isomorphism preserving their ordered distinguished
generators, and also identify a marking with its reverse. Its rank is
the number of generators. The \emph{defect} of a degree-$n$, rank-$r$
marking is $k=n-r$.

Our acquired inputs are the complete transitive rank reduction and
fracture theorem, the theorem on subdivision and contraction at a pure
perfect cut, and the classification by balanced pairs of rooted wings.
Their complete written proofs and accepted finite certificates are
declared dependencies. The first uses the recovered L047 and L044 rank
classifications. The latter two are the new proofs used in the L019
recovery. In particular, no count or completeness assertion from L042
is used as a premise.

\begin{theorem}\label{main}
Let $1\le k\le4$, $n\ge2k+3$, and let a transitive permutation group
of degree $n$ have a connected string C-marking of rank $n-k$.
Its group is $S_n$. Up to marked isomorphism and reversal, the marking
belongs to exactly one of the families in Table~\ref{families}.
There are respectively $1,1,7,9$ families for $k=1,2,3,4$.
All these families exist for every degree in the stated range.

For $n\ge7$ this is also the classification of the given natural
permutation markings up to point conjugacy and reversal. For $n=5,6$
the only case is $k=1$ and the marking is the simplex up to marked
group isomorphism; in degree six one chooses its ordinary action
to obtain the displayed graph. In degrees $n\ge7$ every maximal
omission of a marking in this range is intransitive.
\end{theorem}

We give precise conventions for the table so that every edge of every
source diagram is specified. Put $m=n-k-1$. A word such as
\[
                0,1,\{0,2\},1,2,\ldots,m
\]
denotes a path with the indicated edge labels; a set in a position means
parallel edges on that pair of vertices. Consecutive ascending ranges
include their two endpoints. Thus the word just displayed has initial
entries $0,1,\{0,2\}$ followed by the range $1,\ldots,m$.
The following five additional graphs have vertices in a top row,
numbered from left to right starting at zero, and a short lower row.
Every edge not specified below is absent:
\begin{itemize}
\item $A_m$: top word $2,1,2,3,\ldots,m$; two lower vertices
joined by label $2$; vertical label-$0$ edges from top vertices $0,1$
to the corresponding lower vertices.
\item $G_m$: top word $0,1,\ldots,m$; two lower vertices joined
by label $0$; two parallel vertical edges of labels $1,2$ at top
vertex $0$, and a vertical edge of label $2$ at top vertex $1$.
\item $B_m$: top word $0,1,\ldots,m$; three lower vertices with
word $0,1$; vertical label-$3$ edges at top vertices $0,1,2$.
\item $E_m$: top word $0,1,\ldots,m,m-1$; two lower vertices
joined by label $0$; vertical label-$2$ edges at top vertices $0,1$.
\item $I_m$: the graph $E_m$ with one further vertical edge,
of label $1$, at top vertex $0$.
\end{itemize}
These definitions have exactly $n$ vertices at the defects assigned
in the table. They also resolve the ordinary convention for shortening
the dotted paths at the smallest permitted degree. Our documentary
adapter parses every arrow and label of the original Figure~1 at
degree $31$, expands its dotted paths, and checks literal point
conjugacies to these definitions. The original native figure has
also been inspected directly.

\begin{longtable}{cllc}
\caption{The source families, with a pure cut used for subdivision.}
\label{families}\\
\toprule
$k$ & Source family & Word or graph & Cut $j$\\
\midrule\endfirsthead
\toprule $k$ & Source family & Word or graph & Cut $j$\\\midrule\endhead
1 & simplex & $0,1,\ldots,m$ & 1\\
2 & rank $n-2$ & $1,0,1,2,\ldots,m$ & 2\\
3 & $A$ & $A_m$ & 3\\
3 & $G$ & $G_m$ & 3\\
3 & $K$ & $1,0,1,2,\ldots,m,m-1$ & 2\\
3 & $L_2$ & $2,1,0,1,2,\ldots,m$ & 3\\
3 & $L_1$ & $0,1,0,1,2,\ldots,m$ & 2\\
3 & $M_2$, $i=0$ & $0,1,\{0,2\},1,2,\ldots,m$ & 3\\
3 & $L_3$ & $\{0,2\},1,0,1,2,\ldots,m$ & 3\\
4 & $B$ & $B_m$ & 4\\
4 & $E$ & $E_m$ & 3\\
4 & $I_2$ & $I_m$ & 3\\
4 & $N_1$ & $0,1,0,1,2,\ldots,m,m-1$ & 2\\
4 & $N_2$ & $2,1,0,1,2,\ldots,m,m-1$ & 3\\
4 & $N_3$ & $\{0,2\},1,0,1,2,\ldots,m,m-1$ & 3\\
4 & $O_2$, $i=0$ & $0,1,\{0,2\},1,2,\ldots,m,m-1$ & 3\\
4 & $P_1$ & $1,0,1,0,1,2,\ldots,m$ & 2\\
4 & $Q_2$ & $3,2,1,0,1,2,3,\ldots,m$ & 4\\
\bottomrule
\end{longtable}

The names refer to Table~2 and Section~4 of the original article.
We write $I_2$ following Proposition~4.15; the corresponding graph
in Table~2 and the concluding list is labelled simply $(I)$. The first nine rows of our table
occur in the left column of Figure~1, and the last nine in its right
column, in the same vertical order.

\begin{lemma}\label{growth}
For each row of Table~\ref{families}, its edge with label $j$ is a
pure perfect cut at every $n\ge2k+3$. Subdividing this edge by a
new vertex, replacing its label by the consecutive labels $j,j+1$,
and increasing every former label greater than $j$ by one gives
the same family in degree $n+1$, up to ordered point conjugacy.
\end{lemma}
\begin{proof}
The edge lies in the ascending path between the fixed left part and
the right end. Its label occurs exactly once. Removing it separates
the graph into two connected parts. All edges with smaller labels
lie in the left part, and all edges with larger labels in the right
part; parallel edges and the short lower row occur wholly on the
appropriate side. These statements follow directly from the table.
For completeness, the only right parts with a repeated label have
terminal word $m-1,m,m-1$. At the smallest permitted degree, each
chosen $j$ is strictly less than $m-1$; the inequalities persist as
$n$ increases. Every other repeated or vertical label is smaller
than $j$. Thus neither exception crosses the chosen edge.

The omission has exactly the two connected parts as its orbits,
because orbits of permutation generators are precisely the connected
components of their labelled graph. The smaller and larger marks
act on the respective parts, proving that the cut is perfect.
Its generator is the single transposition on that edge, so it is pure.
Subdividing adds one entry to the ascending range. The left part
is unchanged, the labels of the right part increase by one, and
$m$ changes to $m+1$. These are exactly the defining formulas in
the table, including the terminal repeated label. Map the old
vertices in path order and the lower vertices to their corresponding
vertices; the inserted vertex fills the new path position. This
is an explicit ordered conjugacy for every $n$, not an inference
from finitely many numerical checks.
\end{proof}

\begin{lemma}\label{bases}
At degree $n=2k+3$, the rows of Table~\ref{families} give precisely
the complete classes furnished by the balanced-pair classification,
for each $k=1,2,3,4$.
\end{lemma}
\begin{proof}
We acquire the complete rooted-wing census and its independent
review, rather than repeating that enumeration. A rank-$p$ wing is
an actual transitive fractured C-marking on $2p$ points, rooted
at a point fixed by every mark except the first, and satisfying
the full kernels for adjoining a pure leaf at that root. The
accepted census has $1,1,6,3$ rooted point-conjugacy classes for
$p=1,2,3,4$. Write them $W_{p,s}$, with $s$ one plus their
zero-based index in the retained literal census. Point conjugacy,
not abstract automorphism of a six-point group, is the convention.

The balanced-pair theorem gives one class for every unordered
pair of these wings with ranks summing to $k+1$. It provides literal
full markings obtained by putting the two roots on opposite sides
of a new middle vertex. The certificate in this application records
an explicit point bijection, and whether the target marking is
reversed, between each smallest source diagram and the following
pair:
\[
\begin{array}{c|l}
k=1 & \text{simplex}: (W_{1,1},W_{1,1})\\
k=2 & \text{rank }n-2: (W_{1,1},W_{2,1})\\\hline
k=3 & A:(W_{1,1},W_{3,5}),\quad G:(W_{1,1},W_{3,6}),\\
&K:(W_{2,1},W_{2,1}),\quad L_2:(W_{1,1},W_{3,2}),\\
&L_1:(W_{1,1},W_{3,1}),\quad M_2(0):(W_{1,1},W_{3,4}),\\
&L_3:(W_{1,1},W_{3,3})\\\hline
k=4 & B:(W_{1,1},W_{4,3}),\quad E:(W_{2,1},W_{3,5}),\\
&I_2:(W_{2,1},W_{3,6}),\quad N_1:(W_{2,1},W_{3,1}),\\
&N_2:(W_{2,1},W_{3,2}),\quad N_3:(W_{2,1},W_{3,3}),\\
&O_2(0):(W_{2,1},W_{3,4}),\quad P_1:(W_{1,1},W_{4,1}),\\
&Q_2:(W_{1,1},W_{4,2}).
\end{array}
\]
Each bijection is verified against every point image of every mark.
Its discovery and verification require only the literal permutations:
assign the image of one point, propagate under the ordered generators,
and reject conflicts. Since the action is transitive, every possible
conjugacy is obtained by trying all images of that point. Reversing
the target covers duality. Each row matches exactly one acquired
class, and for each $k$ every class occurs exactly once. This proves
the statement, including existence, IP and distinctness, by the
already proved properties of those complete acquired classes.
\end{proof}

\begin{proof}[Proof of Theorem~\ref{main}]
The threshold implies
\[
                 2r=2(n-k)\ge n+3>n+2.
\]
The acquired transitive fracture theorem therefore makes the group
the natural $S_n$ and all maximal omissions intransitive for $n\ge7$.
It also supplies the simplex conclusion and the stated choice of
action for degrees five and six. Its proof contains the exact small
exceptions to the rank bounds; no generic half-degree estimate is
being substituted for them here.

At $n=2k+3$, Lemma~\ref{bases} gives all classes and their stated
numbers. The acquired stability bijection identifies the classes
in each subsequent degree with those in the preceding degree by
subdivision at an internal pure perfect cut. It also proves that
using any internal pure cut gives the same class, including under
reversal. Lemma~\ref{growth} identifies this operation with growth
inside each displayed family. Induction on $n$ now proves existence,
exhaustion and pairwise distinctness of the table. For $n\ge7$, every
automorphism of $S_n$ is inner, so an ordered marked isomorphism is
point conjugacy. The explicit degree-six qualification completes
the proof in the only exceptional action degree.
\end{proof}

\paragraph{Source contract and qualifications.}
Theorem~\ref{main} recovers Theorem~1.1, Figure~1, the two principal
counts and the symmetric-group conclusion of the abstract. It also
recovers the intransitive-omission conclusion of Corollary~5.2 for
$n\ge7$, with the necessary ordinary-action choice when $n=6$.
In degree six an outer image of the simplex is a different natural
permutation tuple with transitive endpoint omissions; the acquired
small census already exhibits it. Thus a literal statement about
every degree-six action would be false, although the classification
up to marked group isomorphism is unchanged. The original theorem's
printed $1\ge i\ge4$ is read as $1\le i\le4$, consistently with
the abstract and subsequent propositions.

The other cells of the introductory Table~1 are historical ancillary
census reports and are not asserted as additional principal results
of this recovery. The 38 candidate families and their elimination
form the original proof route; we replace that route by the complete
balanced-pair theorem and the explicit identifications above. We do
not claim an independent recovery of every intermediate lemma about
arbitrary candidate graphs. The introductory higher-defect conjectural
direction is covered by the attributed L019 stability result and
is not presented as a newly solved open conjecture here.

The new contribution in this application is the exact conversion from
rooted wings and pure subdivision to all eighteen source diagrams.
The generic theory, finite census and their original attributions
are reused through their accepted editions. No compression figure,
new Lean result or separate credit for a further original is added.
\end{document}
